% Generated draft: requires mathematical and source review.
\documentclass[11pt]{article}
\usepackage[margin=25mm]{geometry}
\usepackage{fontspec}
\setmainfont{Latin Modern Roman}
\usepackage{amsmath,amssymb,mathtools,unicode-math}
\setmathfont{Latin Modern Math}
\usepackage{enumitem,needspace,xurl,hyperref}
\setlength{\parindent}{0pt}
\setlength{\parskip}{5pt}
\setlength{\emergencystretch}{3em}
\newcommand{\N}{\mathbb{N}}
\newcommand{\Z}{\mathbb{Z}}
\newcommand{\Q}{\mathbb{Q}}
\newcommand{\R}{\mathbb{R}}
\newcommand{\C}{\mathbb{C}}
\newcommand{\F}{\mathbb{F}}
\newcommand{\A}{\mathbb{A}}
\newcommand{\PP}{\mathbb{P}}
\newcommand{\E}{\mathbb{E}}
\newcommand{\eps}{\varepsilon}
\newcommand{\dd}{\,\mathrm{d}}
\begin{document}
\section*{UnsolvedMath ID: 700877}
\textbf{Source problem code:} catalogue entry
\subsection{Definitions and mathematical statement}
For the parameter pairs \((s,t)=(4,11)\) and \((s,t)=(5,7)\), consider an incidence system \(\mathcal G=(P,\mathcal L,I)\), where \(P\) is a set of points, \(\mathcal L\) is a set of lines, and \(I\subseteq P\times\mathcal L\) is the incidence relation. Write
\[
P_L=\{p\in P:(p,L)\in I\},\qquad
\mathcal L_p=\{L\in\mathcal L:(p,L)\in I\}.
\]
The supplied definition of a generalized quadrangle \(GQ(s,t)\) requires
\[
\begin{aligned}
&|P_L|=s+1 &&\text{for every }L\in\mathcal L,\\
&|P_L\cap P_M|\leq 1 &&\text{for all distinct }L,M\in\mathcal L,\\
&|\mathcal L_p|=t+1 &&\text{for every }p\in P,
\end{aligned}
\]
and, for every \(p\in P\) and \(L\in\mathcal L\) with \(p\notin P_L\),
\[
\left|\{M\in\mathcal L_p:P_M\cap P_L\neq\varnothing\}\right|=1.
\]
For the intended questions below, take \(P\neq\varnothing\). This nonemptiness requirement makes explicit the conventional intended reading; the supplied definition does not state it, and its listed axioms alone permit an empty incidence system. The record also does not explicitly impose finiteness on \(P\) or \(\mathcal L\); no separate finiteness hypothesis is added here.

An incidence-preserving automorphism of \(\mathcal G\) is a pair of bijections \(\sigma_P:P\to P\) and \(\sigma_{\mathcal L}:\mathcal L\to\mathcal L\) such that
\[
(p,L)\in I\quad\Longleftrightarrow\quad
(\sigma_P(p),\sigma_{\mathcal L}(L))\in I.
\]
Write \(\operatorname{Aut}(\mathcal G)\) for this group. An involution is an element \(\sigma\in\operatorname{Aut}(\mathcal G)\) satisfying \(\sigma^2=1\) and \(\sigma\neq 1\).

\begin{enumerate}
\item[(a)] Does there exist a nonempty generalized quadrangle \(GQ(4,11)\) satisfying the stated axioms?
\item[(b)] For a hypothetical nonempty generalized quadrangle \(\mathcal G\) with parameters \((5,7)\), can \(\operatorname{Aut}(\mathcal G)\) contain an involution? Equivalently, can there exist such a pair \((\mathcal G,\sigma)\) with \(\sigma\in\operatorname{Aut}(\mathcal G)\), \(\sigma\neq 1\), and \(\sigma^2=1\)? This question does not assume or assert the existence of a \(GQ(5,7)\).
\end{enumerate}
Here automorphisms preserve points and lines separately, as in the usual incidence-automorphism convention; the supplied record does not explicitly discuss incidence symmetries interchanging these types.
\subsection{Short English statement}
The questions ask whether a generalized quadrangle with parameters \((4,11)\) exists and whether a hypothetical generalized quadrangle with parameters \((5,7)\) can have an automorphism of order two. Nonemptiness and the convention that automorphisms preserve points and lines separately are made explicit and require source review.
\subsection{Sources}
\begin{itemize}
\item Local UnsolvedMath registry, permanent ID \(700877\), exact source problem code \texttt{catalogue entry}, title \emph{Kourovka Notebook Problem 12.58}. The supplied record labels the problem \texttt{open}; this is a supplied registry classification, not an independently verified mathematical or literature assessment.
\item \emph{The Kourovka Notebook (Unsolved Problems in Group Theory)}, original problem \(12.58\), 12th issue (1992), attributed to A. A. Makhn\"ev. The supplied background identifies the 21st issue edition (2026), \texttt{arXiv:1401.0300v48}. This bibliographic attribution is supplied evidence and has not been independently checked.
\item The supplied collection description identifies E. I. Khukhro and V. D. Mazurov as editors of the Kourovka Notebook collection covering Issues 1--21 (1965--2026). The registry record was updated on 2026-10-16; this metadata date is not a source-verification date. No separate dated literature assessment is supplied.
\end{itemize}
\end{document}
